\documentclass[11pt]{article}
\usepackage[margin=1in]{geometry}
\usepackage{amsmath,amssymb,amsthm}
\usepackage[hidelinks]{hyperref}

\newtheorem{theorem}{Theorem}
\newtheorem{proposition}[theorem]{Proposition}
\newtheorem{lemma}[theorem]{Lemma}
\newtheorem{corollary}[theorem]{Corollary}
\theoremstyle{definition}
\newtheorem{remark}[theorem]{Remark}

\newcommand{\R}{\mathbb{R}}
\newcommand{\norm}[1]{\left\|#1\right\|}
\newcommand{\abs}[1]{\left|#1\right|}
\renewcommand{\d}{\,d}

\title{Sharp Faber--Krahn stability via the Kohler--Jobin transfer}
\author{Plan 1 / Sharp stability for Faber--Krahn}
\date{2026-05-11}

\begin{document}
\maketitle

\begin{abstract}
Route A gives sharp \emph{Saint--Venant} stability through the single-set
selection map, graph entry, interpolation into the nearly spherical class,
and BDV Theorem~3.3:
\[
E(\Omega)-E(B_1)\ge c_*(N,R)\,\alpha(\Omega)
\quad\hbox{in the small-\(\alpha\) regime}.
\]
Together with BDV Lemma 4.2
($\mathcal A(\Omega)^2\lesssim\alpha(\Omega)$) and the far-from-ball
Saint--Venant bound, this gives the global energy-form estimate
\[
E(\Omega)-E(B^*)\ge c_{\rm SV}(N,R)\mathcal A(\Omega)^2.
\]
The Bernoulli spectral closure of
\texttt{fixed-domain-bernoulli-expansion.md} is a possible Route B, but it is
not used for the unconditional route recorded here.

This note records the Kohler--Jobin transfer that converts this into sharp
\emph{Faber--Krahn} stability for the first Dirichlet eigenvalue:
\[
\lambda_1(\Omega)-\lambda_1(B^*)\ge c_{\rm FK}(N,R)\,\mathcal A(\Omega)^2,
\]
where $B^*$ is the ball with $\abs{B^*}=\abs{\Omega}$. The transfer uses
only the Kohler--Jobin inequality (a non-perturbative pointwise comparison
that holds for every open set) plus elementary scaling. The output is the
BDV sharp Faber--Krahn theorem.
\end{abstract}

\section{The Kohler--Jobin inequality}\label{sec:kj}

For an open bounded set $\Omega\subset\R^N$, let
\[
T(\Omega):=\int_\Omega u_\Omega\d x,\qquad
-\Delta u_\Omega=1\hbox{ in }\Omega,\;u_\Omega=0\hbox{ on }\partial\Omega,
\]
be the torsional rigidity, and let $\lambda_1(\Omega)$ be the first Dirichlet
eigenvalue of $-\Delta$ on $\Omega$. Both functionals scale as
\[
T(t\Omega)=t^{N+2}T(\Omega),\qquad\lambda_1(t\Omega)=t^{-2}\lambda_1(\Omega),
\]
so the product
\[
\Psi(\Omega):=T(\Omega)\cdot\lambda_1(\Omega)^{(N+2)/2}
\]
is scale-invariant.

\begin{theorem}[Kohler--Jobin, 1978]\label{thm:kj}
For every open bounded $\Omega\subset\R^N$ with $\lambda_1(\Omega)<\infty$,
\[
\Psi(\Omega)\ge\Psi(B),
\]
where $B$ is any ball, with equality iff $\Omega$ is a ball (up to translation
and null sets).
\end{theorem}

See Kohler--Jobin~(1978) and Brasco, ``On torsional rigidity and principal
frequencies'' (ESAIM:COCV, 2014) for proofs via rearrangement. The constant
$\Psi(B)$ is universal:
\[
\Psi(B)=T(B)\cdot\lambda_1(B)^{(N+2)/2}=c_N
\]
depends only on $N$ (not on the radius), since $\Psi$ is scale-invariant.

\section{Transfer at fixed volume}\label{sec:transfer}

Fix $\abs{\Omega}=\abs{B^*}$ where $B^*$ is the comparison ball. Setting
$L:=\lambda_1(B^*)$, $T_0:=T(B^*)$, and abbreviating $\beta:=(N+2)/2$,
Theorem~\ref{thm:kj} reads
\begin{equation}\label{eq:kjfixed}
T(\Omega)\cdot\lambda_1(\Omega)^\beta\ge T_0\cdot L^\beta.
\end{equation}

\begin{lemma}[Pointwise transfer]\label{lem:pointwise}
For every $\Omega$ with $\abs{\Omega}=\abs{B^*}$,
\[
\lambda_1(\Omega)\ge L\cdot\Bigl(\frac{T_0}{T(\Omega)}\Bigr)^{1/\beta}.
\]
Equivalently, with the Saint--Venant deficit
$\delta_{\rm SV}(\Omega):=T_0-T(\Omega)\ge0$ (note $\delta_{\rm SV}\ge0$ at fixed volume by
Saint--Venant),
\[
\lambda_1(\Omega)-L\ge L\cdot\Bigl[\Bigl(\frac{1}{1-\delta_{\rm SV}/T_0}\Bigr)^{1/\beta}-1\Bigr].
\]
\end{lemma}

\begin{proof}
Rearrange~\eqref{eq:kjfixed}.
\end{proof}

\begin{lemma}[Small-deficit linearization]\label{lem:linear}
For $\delta_{\rm SV}\le T_0/2$,
\[
\lambda_1(\Omega)-L\ge\frac{L}{\beta T_0}\cdot\delta_{\rm SV}
=\frac{2L}{(N+2)T_0}\cdot\bigl(T_0-T(\Omega)\bigr).
\]
\end{lemma}

\begin{proof}
For $t\in[0,1/2]$,
$(1-t)^{-1/\beta}-1\ge t/\beta$ by elementary convexity. Apply with
$t=\delta_{\rm SV}/T_0$.
\end{proof}

\begin{remark}
The factor $2L/((N+2)T_0)$ is universal in $N$ once normalized. For the unit
ball $B_1$ in $\R^N$, $L=L_N$ is the first Dirichlet eigenvalue of $-\Delta$
on $B_1$ and $T_0=T_N$ is the torsional rigidity of $B_1$. Both are explicit:
$L_N$ is the smallest positive zero of a Bessel function, $T_N=\abs{B_1}/(N(N+2))$.
\end{remark}

\section{Sharp Faber--Krahn from sharp Saint--Venant}

\begin{proposition}[Faber--Krahn transfer from energy Saint--Venant]\label{prop:svtofk}
Assume the sharp Saint--Venant stability holds in energy form
\begin{equation}\label{eq:svsharp}
E(\Omega)-E(B^*)\ge c_{\rm SV}(N,R)\,\mathcal A(\Omega)^2
\quad\hbox{for all }\Omega\subset B_R\hbox{ with }\abs{\Omega}=\abs{B^*}.
\end{equation}
Then for the same class of sets,
\[
\lambda_1(\Omega)-\lambda_1(B^*)
\ge c_{\rm KJ}(N,R)\,\mathcal A(\Omega)^2,
\]
where
\[
c_{\rm KJ}(N,R):=
\min\left\{
\frac{4L\,c_{\rm SV}(N,R)}{(N+2)T_0},
\frac{L(2^{2/(N+2)}-1)}{4}
\right\}.
\]
\end{proposition}

\begin{proof}
Set $\delta_{\rm SV}:=T_0-T(\Omega)=2(E(\Omega)-E(B^*))$. If
$\delta_{\rm SV}\le T_0/2$, Lemma~\ref{lem:linear} and~\eqref{eq:svsharp}
give
\[
\lambda_1(\Omega)-L
\ge \frac{2L}{(N+2)T_0}\delta_{\rm SV}
\ge \frac{4L\,c_{\rm SV}(N,R)}{(N+2)T_0}\mathcal A(\Omega)^2.
\]
If $\delta_{\rm SV}>T_0/2$, Lemma~\ref{lem:pointwise} gives
\[
\lambda_1(\Omega)-L\ge L\left(2^{2/(N+2)}-1\right).
\]
Since $\mathcal A(\Omega)\le2$, this is at least
$L(2^{2/(N+2)}-1)\mathcal A(\Omega)^2/4$.
\end{proof}

\begin{theorem}[Sharp Faber--Krahn stability]\label{thm:fk}
There is $c_{\rm FK}(N,R)>0$ such that, for every open $\Omega\subset B_R$
with $\abs{\Omega}=\abs{B^*}$,
\[
\lambda_1(\Omega)-\lambda_1(B^*)\ge c_{\rm FK}(N,R)\,\mathcal A(\Omega)^2.
\]
\end{theorem}

\begin{proof}
Route A proves the global energy-form Saint--Venant estimate
\[
E(\Omega)-E(B^*)\ge c_{\rm SV}(N,R)\mathcal A(\Omega)^2
\]
after the small-\(\alpha\) argument is combined with the BDV suboptimal
far-from-ball Saint--Venant bound. Applying Proposition~\ref{prop:svtofk}
gives the claim with \(c_{\rm FK}:=c_{\rm KJ}\).
\end{proof}

\section{Quantitative dependence on \texorpdfstring{$(N,R)$}{(N,R)}}

The constant $c_{\rm FK}(N,R)$ is the minimum of the small-deficit and
large-deficit Kohler--Jobin branches:

\begin{itemize}
\item \textbf{Small-deficit constant:}
\[
c_1(N,R)=\frac{4L_N\,c_{\rm SV}(N,R)}{(N+2)T_N},
\]
where \(c_{\rm SV}(N,R)\) is the energy-form Saint--Venant constant from
Route A. Tracing back through the selection chain:
\begin{itemize}
\item Selection (\texttt{quantitative-selection-principle.md}) preserves
the quotient $Q_\alpha$ up to the factor \(2C_{\rm sel}(N,R)\).
\item Graph entry (\texttt{graph-entry-closure.md}) requires
$\alpha(\widetilde\Omega)\le\alpha_{\rm graph}(N,R)$.
\item Interpolation from the small \(L^\infty\) graph norm and uniform
Schauder bounds gives
\(\norm{g}_{C^{2,\gamma_0}}\le\delta_{\rm sph}(N,\gamma_0)\).
\item BDV Theorem~3.3 gives the nearly spherical lower bound.
\end{itemize}
Combining, the small-\(\alpha\) quotient constant is
\(q_*=c_*(N)/(2C_{\rm sel}(N,R))\), and the global energy-form
Saint--Venant constant is
\[
c_{\rm SV}(N,R)=
\min\{q_*/C_{\rm Fr}(N),\,c_{\rm far,SV}(N,R)\}.
\]
Thus
\[
c_1(N,R)=\frac{4L_N\,c_{\rm SV}(N,R)}{(N+2)T_N}.
\]
\item \textbf{Large-deficit constant:}
\[
c_2(N)=\frac{L_N(2^{2/(N+2)}-1)}{4}.
\]
This branch uses only the pointwise Kohler--Jobin inequality and
\(\mathcal A\le2\); it does not require an additional Hausdorff compactness
argument for eigenvalues.
\end{itemize}

\begin{remark}[Comparison with BDV]
This is exactly the chain BDV use to prove sharp Faber--Krahn: their Theorem
1.1 is our Theorem~\ref{thm:fk}. The reliable new content of Plan 1 is the
single-set, constant-tracked replacement for the selection/compactness step.
The final Route A closure still uses BDV's nearly spherical theorem. The
Kohler--Jobin transfer step is shared with BDV and uses no new input.
\end{remark}

\section{Status}

The full chain for sharp Faber--Krahn stability is structurally complete,
with two possible closure routes:
\begin{enumerate}
\item \textbf{Route A (unconditional):} BDV's nearly spherical second
variation supplies the sharp Saint--Venant constant $c_{\rm SV}(N,R)$.
\item \textbf{Route B (conditional):} the Bernoulli spectral closure of
\texttt{fixed-domain-bernoulli-expansion.md}, supplemented by
\texttt{corrections-response.md}, supplies $c_{\rm SV}(N,R)$, conditional on
the explicit constants in the second-variation source enumeration.
\end{enumerate}
Either way, the Kohler--Jobin transfer of Lemma~\ref{lem:pointwise} +
Lemma~\ref{lem:linear} + the large-deficit branch (Theorem~\ref{thm:fk}) converts
$c_{\rm SV}$ into $c_{\rm FK}$.

The quantitative Faber--Krahn constant $c_{\rm FK}(N,R)$ is given by an
explicit formula in terms of:
\begin{itemize}
\item the universal ball constants $L_N=\lambda_1(B_1)$ and
$T_N=T(B_1)=\abs{B_1}/(N(N+2))$;
\item the BDV regularity constants from selection and graph entry;
\item the interpolation threshold and the BDV nearly spherical constants
\(c_{\rm sph},\delta_{\rm sph}\).
\end{itemize}

The exponent is sharp: the inequality
$\lambda_1(\Omega)-\lambda_1(B^*)\ge c_{\rm FK}\mathcal A(\Omega)^2$ is the
optimal power of $\mathcal A$, as already observed by Bhattacharya, Hansen,
Nadirashvili (and saturated by ellipsoidal perturbations of the ball).

\end{document}
